\documentclass[tikz,border=5pt,12pt]{standalone}
\usepackage{resistencia}
\input{datos-e01}
\pgfkeys{/pgf/number format/.cd,fixed,precision=1,set decimal separator={,},zerofill=false}
\begin{document}
\begin{tikzpicture}[x=1cm,y=1cm,etq/.style={etiq,fill=white,inner sep=1.2pt}]
\pgfmathsetmacro{\k}{0.075}% cm por MPa
\pgfmathsetmacro{\xuno}{\Suno*\k}\pgfmathsetmacro{\xdos}{\Sdos*\k}\pgfmathsetmacro{\xtres}{\Stres*\k}
\pgfmathsetmacro{\cA}{(\xuno+\xtres)/2}\pgfmathsetmacro{\rA}{(\xuno-\xtres)/2}
\pgfmathsetmacro{\cB}{(\xuno+\xdos)/2}\pgfmathsetmacro{\rB}{(\xuno-\xdos)/2}
\pgfmathsetmacro{\cC}{(\xdos+\xtres)/2}\pgfmathsetmacro{\rC}{(\xdos-\xtres)/2}
\pgfmathsetmacro{\px}{\SigmaPlano*\k}\pgfmathsetmacro{\py}{\TauPlano*\k}
% región admisible: dentro del círculo grande y fuera de los dos pequeños (τ ≥ 0)
\begin{scope}
  \clip (\xtres-0.2,0) rectangle (\xuno+0.2,\rA+0.2);
  \fill[even odd rule,fill=acero!18] (\cA,0) circle (\rA) (\cB,0) circle (\rB) (\cC,0) circle (\rC);
\end{scope}
% ejes
\draw[eje] (\xtres-0.9,0) -- (\xuno+1.1,0) node[right,etq]{$\sigma$ (MPa)};
\draw[eje] (0,-0.2) -- (0,\rA+0.8) node[above,etq]{$\tau$ (MPa)};
% semicírculos
\draw[acero,very thick] (\xuno,0) arc[start angle=0,end angle=180,radius=\rA];
\draw[verde,thick] (\xuno,0) arc[start angle=0,end angle=180,radius=\rB];
\draw[cobre,thick] (\xdos,0) arc[start angle=0,end angle=180,radius=\rC];
% tensiones principales
\foreach \x/\n/\v in {\xtres/3/\Stres, \xdos/2/\Sdos, \xuno/1/\Suno}{
  \draw (\x,0.07) -- (\x,-0.07);
  \node[etq,below=2pt] at (\x,0) {$\sigma_\n=\pgfmathprintnumber{\v}$};
}
% tau maxima
\draw[dashed,tinta] (0,\rA) -- (\cA,\rA);
\pgfmathsetmacro{\taumax}{\rA/\k}
\node[etq,left] at (0,\rA) {$\tau_{\max}=\pgfmathprintnumber{\taumax}$};
% punto de un plano concreto
\draw[dotted,tinta] (\px,0) -- (\px,\py) -- (0,\py);
\fill[violeta] (\px,\py) circle (2.4pt);
\node[etq,text=violeta,above right=1pt] at (\px,\py) {$(\sigma_n,\tau)$};
\end{tikzpicture}
\end{document}
